\documentclass{article}
\usepackage{amsmath, amssymb, amsthm}

\setlength{\parindent}{0pt}

\newtheorem*{definition}{Definition}
\newtheorem*{claim}{Claim}

\begin{document}

\begin{definition}
    A real number $r$ is called sensible if there exist positive integers $a$ and $b$
    such that $\sqrt{a/b} = r$.
\end{definition}

\begin{claim}
    $\sqrt[3]{2}$ is not sensible.
\end{claim}

\begin{proof}
    Suppose for contradiction that $\sqrt[3]{2}$ is sensible. Then there exist positive 
    integers $a$ and $b$ such that $\sqrt{a/b} = \sqrt[3]{2}$. Squaring both sides gives:
    \[
    \frac{a}{b} = \sqrt[3]{4}
    \]
    Since $a/b$ is rational, it follows that $\sqrt[3]{4}$ is rational.

    \medskip

    Since $\sqrt[3]{4}$ is rational, we can write it as a fraction $x/y$ in lowest terms,
    where $x$ is an integer and $y$ is a positive integer. Therefore, we have:

    \[
    \begin{aligned}
    \sqrt[3]{4} &= \frac{x}{y} \\
    4 &= \frac{x^3}{y^3} \\
    4y^3 &= x^3
    \end{aligned}
    \]

    In the last equation, the left side is even as the coefficient is a multiple of 2.
    Thus the right side must also be even. Since $x^3$ is even, $x$ itself must be even.
    Since $x$ is even, we can write $x = 2k$ for some $k \in \mathbb{Z}$. Hence:
    \[
    x^3 = (2k)^3 = 8k^3
    \]
    Substituting into $4y^3 = x^3$ gives
    \[
    4y^3 = 8k^3,
    \]
    so $y^3 = 2k^3$, thus $y^3$ is even. Since $y^3$ is even, $y$ itself must be even.

    \medskip

    However, $x$ and $y$ are both even, contradicting that $x/y$ is in lowest terms.
    Thus $\sqrt[3]{4}$ is irrational, contradicting the assumption that $\sqrt[3]{4}$ 
    is rational.

    \medskip

    Hence $\sqrt[3]{2}$ is not sensible.
\end{proof}

\end{document}
